% =====================================================================
%  1bat-ud3-6.tex  —  HORA 6: Repartiment de la pobresa: proporcionalitat inversa
%  (sense preàmbul: s'inclou des de main.tex)
% =====================================================================
\horatitol{Sessió 6 — Repartiment de la pobresa: proporcionalitat inversa}%
{Estudi de cas (Activitat 2): un fons repartit entre famílies, la hipèrbola i la idea d'asímptota, amb doble lectura matemàtica i social.}

\subsection{Idea clau}

Fins ara hem treballat funcions lineals (rectes). Avui apareix un tipus nou de
relació amb un fort sentit social: la \textbf{proporcionalitat inversa}. El cas és
real (Activitat 2): un \textbf{fons d'emergència} de $\num{120000}\eur$ s'ha de
repartir \textbf{a parts iguals} entre les famílies afectades per un desnonament. El
que rep cada família depèn de \emph{quantes} famílies se'l reparteixin.

\begin{idea}
\textbf{Proporcionalitat inversa:} $\quad f(x)=\dfrac{k}{x}$
\begin{itemize}[leftmargin=1.4em, itemsep=2pt]
  \item Com més gran és $x$, més \textbf{petit} és $f(x)$: el producte
        $x\cdot f(x)=k$ es manté \textbf{constant} (aquí $k=\num{120000}$).
  \item La gràfica és una \textbf{hipèrbola}: decreix de pressa i després cada cop
        més a poc a poc.
  \item \textbf{Asímptota:} quan $x$ es fa molt gran, $f(x)$ s'acosta a $0$ \emph{sense
        arribar-hi mai}. La recta $y=0$ (l'eix horitzontal) és una \textbf{asímptota}.
\end{itemize}
\textbf{Domini lògic:} no es pot repartir entre $0$ famílies ($x\neq 0$) i el nombre
de famílies ha de ser un \textbf{enter positiu}: $\mathrm{Dom}\,f=\{1,2,3,\dots\}$.
\end{idea}

\subsection{Exemples}

\begin{exemple}[el fons de l'Activitat 2]
Amb $f(x)=\dfrac{\num{120000}}{x}$, calculem què rep cada família segons quantes
n'hi hagi:
\begin{center}
\begin{tabular}{c c c c c c}
\toprule
Famílies $x$ & $1$ & $2$ & $4$ & $10$ & $100$ \\
\midrule
Per família $f(x)$ ($\eur$) & $\num{120000}$ & $\num{60000}$ & $\num{30000}$ & $\num{12000}$ & $\num{1200}$ \\
\bottomrule
\end{tabular}
\end{center}
A mesura que augmenten les famílies, l'ajut per família \textbf{cau de pressa} al
principi i després cada cop més lentament. Mai no arriba a $0$: sempre queda alguna
cosa, però pot ser tan petita que pràcticament desaparegui.
\end{exemple}

\begin{exemple}[la forma de la hipèrbola i la seva asímptota]
Aquesta és la gràfica que veureu a GeoGebra (en escala reduïda). Fixeu-vos com la
corba s'acosta a l'eix horitzontal sense tocar-lo:

\begin{center}
\begin{tikzpicture}
\begin{axis}[udaxis, width=9.4cm, height=5.6cm,
    xlabel={\small famílies ($x$)}, ylabel={\small \euro\ per família (en milers)},
    xmin=0, xmax=12, ymin=0, ymax=70,
    xtick={0,2,4,6,8,10}, ytick={0,20,40,60}]
  \addplot[udblue, domain=1.8:12, samples=120]{120/x};
  \addplot[udgrey, dashed, domain=0:12, samples=2]{0};
  \node[udgrey, font=\scriptsize, anchor=south] at (axis cs:6,2) {asímptota $y=0$};
\end{axis}
\end{tikzpicture}

\figcap{$f(x)=\num{120000}/x$: la corba decreix i s'acosta a l'eix horitzontal (asímptota) sense arribar-hi.}
\end{center}
\end{exemple}

\subsection{Exercicis resolts}

\begin{exresolt}[domini lògic i lectura de valors]
Per al fons $f(x)=\dfrac{\num{120000}}{x}$:
\begin{enumerate}[label=\alph*), leftmargin=1.6em]
  \item Per què $x$ no pot valer $0$? I per què no pot ser $2,5$?
  \item Quant rep cada família si n'hi ha $6$? I si n'hi ha $50$?
\end{enumerate}
\solucio
\textbf{a)} No es pot dividir entre $0$: repartir un fons «entre cap família» no té
sentit (ni matemàtic ni real). I $x=2,5$ tampoc val, perquè no hi pot haver «dues
famílies i mitja»: $x$ ha de ser un \textbf{enter positiu}.

\textbf{b)} Substituïm:
\[
f(6)=\frac{\num{120000}}{6}=\num{20000}\eur,\qquad
f(50)=\frac{\num{120000}}{50}=\num{2400}\eur.
\]
\resposta{a) $x\neq 0$ i $x$ enter positiu. \quad b) $\num{20000}\eur$ amb $6$ famílies; $\num{2400}\eur$ amb $50$.}
\end{exresolt}

\begin{exresolt}[interpretar l'asímptota socialment]
Què passa amb l'ajut per família si el nombre de famílies es fa \emph{molt gran}
(per exemple, $\num{1000}$ o $\num{5000}$ famílies)? Quina lectura social en fas?
\solucio
Calculem-ne un parell:
\[
f(\num{1000})=\frac{\num{120000}}{\num{1000}}=120\eur,\qquad
f(\num{5000})=\frac{\num{120000}}{\num{5000}}=24\eur.
\]
Com més famílies, menys rep cadascuna: l'ajut s'acosta a $0$ sense arribar-hi
(l'asímptota $y=0$). \textbf{Lectura social:} un fons fix, per gran que sigui, es
dilueix si s'ha de repartir entre moltíssima gent; arriba un punt en què l'ajut per
família és tan petit que deixa de ser útil. És una metàfora de l'\textbf{esgotament
dels recursos} quan la necessitat creix.
\resposta{L'ajut tendeix a $0$ (asímptota); socialment, el fons es dilueix i deixa de
fer efecte si hi ha massa famílies a atendre.}
\end{exresolt}

\subsection{Exercicis per fer}

\noindent\textit{Pots fer les gràfiques a GeoGebra escrivint, per exemple,
\texttt{f(x)=120000/x} i ajustant la finestra.}

\begin{exercicis}
\begin{exlist}
  \item Completa la taula de valors de $f(x)=\dfrac{\num{120000}}{x}$:
        \begin{center}
        \begin{tabular}{c c c c c}
        \toprule
        $x$ & $3$ & $8$ & $20$ & $200$ \\
        \midrule
        $f(x)$ ($\eur$) & \rule{1.6cm}{0.4pt} & \rule{1.6cm}{0.4pt} & \rule{1.6cm}{0.4pt} & \rule{1.6cm}{0.4pt} \\
        \bottomrule
        \end{tabular}
        \end{center}
  \item Un altre fons, més petit, és $g(x)=\dfrac{\num{6000}}{x}$, on $x$ són les
        persones beneficiàries. Calcula $g(1)$, $g(5)$, $g(20)$ i $g(100)$, i digues
        què li passa a l'ajut quan augmenten les persones.
  \item Indica el \textbf{domini lògic} de cada funció segons el seu context:
        \begin{enumerate}[label=\alph*), leftmargin=1.6em, topsep=2pt]
          \item $f(x)=\dfrac{\num{120000}}{x}$, on $x$ és el nombre de famílies.
          \item $h(x)=\dfrac{900}{x}$, on $x$ és el nombre de torns de voluntariat per
                repartir $900$ àpats.
        \end{enumerate}
  \item Dibuixa a GeoGebra $f(x)=\dfrac{\num{120000}}{x}$ i respon: cap a quin valor
        s'acosta la corba quan $x$ es fa molt gran? Com es diu aquesta recta a la qual
        s'acosta?
  \item \textbf{(Comparació)} Tens dos fons: $A(x)=\dfrac{\num{120000}}{x}$ i
        $B(x)=\dfrac{\num{60000}}{x}$. Amb el mateix nombre de famílies, quin reparteix
        més per família? Per què? Què tenen en comú les dues gràfiques quan $x$ és molt
        gran?
  \item \textbf{(Raonament social)} Algú proposa: «com que sempre queda alguna cosa
        per repartir, podem atendre tantes famílies com vulguem». Explica per què
        aquest raonament és matemàticament cert però socialment problemàtic, fent
        servir la idea d'asímptota.
\end{exlist}
\end{exercicis}

% ---------------------------------------------------------------------
\fullsolucions{Sessió 6}
\begin{solucions}
\begin{sollist}
  \item $f(3)=\num{40000}\eur$; \quad $f(8)=\num{15000}\eur$; \quad
        $f(20)=\num{6000}\eur$; \quad $f(200)=600\eur$.
  \item $g(1)=\num{6000}\eur$, $g(5)=\num{1200}\eur$, $g(20)=300\eur$,
        $g(100)=60\eur$. Quan augmenten les persones, l'ajut per persona decreix i
        s'acosta a $0$.
  \item (a) $\mathrm{Dom}=\{1,2,3,\dots\}$ (enter positiu de famílies). \quad
        (b) $\mathrm{Dom}=\{1,2,3,\dots\}$ (enter positiu de torns; a més, ha de ser
        un divisor raonable de $900$ si es vol repartiment exacte).
  \item S'acosta a $y=0$ (l'eix horitzontal); aquesta recta és l'\textbf{asímptota}
        horitzontal.
  \item Amb el mateix $x$, el fons $A$ reparteix \emph{el doble} per família, perquè
        té el doble de diners ($\num{120000}$ contra $\num{60000}$). Totes dues
        gràfiques, però, s'acosten a $0$ (la mateixa asímptota $y=0$) quan $x$ és molt
        gran.
  \item Matemàticament $f(x)$ mai no és exactament $0$, per tant «sempre queda alguna
        cosa». Però socialment, quan l'ajut per família s'acosta a $0$, esdevé tan
        petit que ja no cobreix cap necessitat real: repartir entre massa famílies
        equival, a la pràctica, a no ajudar ningú. L'asímptota mostra el límit real
        del fons.
\end{sollist}
\end{solucions}
